\documentclass[10pt,a4paper]{article}
\usepackage[margin=1.55cm]{geometry}
\usepackage{xcolor}
\usepackage{amsmath,amssymb}
\usepackage{graphicx}
\usepackage{enumitem}
\usepackage{fontspec}
\usepackage[CJKmath=true]{xeCJK}
% CJK main font is resolved at COMPILE time on the actual machine, not baked in at
% generation time, so a .tex generated on one OS still renders on another (the older
% Python-side probe hardcoded one OS's font name into the file and broke when moved).
% Walk a cross-platform priority list and use the first font that exists; AutoFakeBold
% + AutoFakeSlant give bold/italic CJK even on faces that ship no bold/italic cut.
\newcommand{\setCJKto}[1]{\setCJKmainfont{#1}[AutoFakeBold=true,AutoFakeSlant=0.2]}
\IfFontExistsTF{Noto Sans CJK SC}{\setCJKto{Noto Sans CJK SC}}{%        % Linux (fonts-noto-cjk); cross-platform if installed
 \IfFontExistsTF{Source Han Sans SC}{\setCJKto{Source Han Sans SC}}{%   % Adobe 思源黑体; cross-platform
  \IfFontExistsTF{PingFang SC}{\setCJKto{PingFang SC}}{%               % macOS 10.11+
   \IfFontExistsTF{Hiragino Sans GB}{\setCJKto{Hiragino Sans GB}}{%    % older macOS
    \IfFontExistsTF{Microsoft YaHei}{\setCJKto{Microsoft YaHei}}{%     % Windows
     \IfFontExistsTF{WenQuanYi Micro Hei}{\setCJKto{WenQuanYi Micro Hei}}{% % older Linux
      \setCJKto{Songti SC}}}}}}}                                       % macOS bottom fallback
\definecolor{cblue}{HTML}{4C72B0}
\setlength{\parskip}{3pt}\setlength{\parindent}{0pt}
\setlist[enumerate]{topsep=2pt,itemsep=2pt,parsep=1pt,partopsep=0pt}
\setlist[itemize]{topsep=2pt,itemsep=2pt,parsep=1pt,partopsep=0pt}
% Let prose-embedded inline math break across lines and absorb small overflows, so simple
% formulas inserted mid-sentence stop running past the right margin.
\tolerance=1200
\emergencystretch=2.5em
\relpenalty=20
\binoppenalty=40
\hfuzz=1pt
\begin{document}

\section*{生成测试的断言检查了需求中的哪些约束？}
\textbf{方法名称：} Requirement-guided assertion replay
\section*{研究动机}
Req2Test-Bench 的 8 个项目中，我们复核了 44 条已经通过的测试。8 条（18.2\%）漏查了部分业务要求，分布在 7 个项目。其中 4 条的问题出在断言，也就是决定结果能否通过的那段代码。一个 python-docx 测试检出了官方用于计分的全部 93 个变异体。把原 Word 文档的段落放在原位，取出三张内容不同的表格，按 6 种顺序分别交给原测试。它全都放行，只有 1 种符合文档顺序。这次用本地 python-docx 读出了真实表格，现有记录还不能确认发布版本曾返回错序。

搜索文件恰好说明了另一种要求。ast-grep 可以先列 a.js，也可以先列 b.js，但两个匹配文件各应出现一次。原测试接受两种合法顺序，也接受连续两次列出 a.js。若直接把顺序写死，又会错拒一个合法结果。需求和样例文件足以判断这些具体结果，再把它们交给原断言，就能看到它究竟判错了哪一项。

TestPrism 已经可以在不同实现上评价整套测试，Ma 和 Eisty 也已从文字需求直接生成判定代码。这里要回答的问题是，拿到一条原测试确实判错的结果，能否比直接照着需求改断言多查出独立缺陷，同时让原先能通过的合法变体继续通过。这个效果仍待验证。
\section*{方法}
\subsection*{背景}
\begin{enumerate}[leftmargin=*,start=1]
  \item 逐项目记录需求原句、测试执行的动作，以及接收结果的代码行。筛过多少条、多少条没触发动作或看不到结果，都要计数。TestSuiteBench 的 165 份规格中有 14 份明确允许输出换序，可从中找案例。调试前预留新仓库。已有 3 次重放和 44 条人工复核的通过样本用于开发，这 44 条已涉及 Req2Test-Bench 全部 8 个项目。同一测试、同一条款和同一份样例记一行。保存版本，以及执行动作、接收结果、作出判断的三处代码位置。同一任务若有多条需求，分别统计任务数和条款数。
\par\emph{为什么：} 测试先要执行对应动作、读到相关结果，改断言才有机会补上检查。后续来自同一项目的样本可以检验结果能否在该项目内重现。
\end{enumerate}
\subsection*{按需求构造结果，重放原断言}
\emph{找出一份需求与原测试判定不一致的具体结果。}
\begin{enumerate}[leftmargin=*,start=2]
  \item 固定模型版本、提示词和生成设置，仅提供公开需求、原测试、样例及允许查看的接口说明。核对适用于当前场景的原句，查清能靠什么把返回元素和样例中的对象对应起来。按条款记录每个元素应出现几次、完整顺序，或谁必须在谁之前。按已有依据造一份类型正确的合规结果。若多项都能排在下一位，选前置条件已满足且在样例中最靠前的。规则不明、次序矛盾或对象无法辨认，均记为未知。Word 案例生成时已打印出中间表独有的文字，后来复核原件又从结构上区分了首尾两张表。
\par\emph{为什么：} “每个文件出现一次”“按文档顺序”和“父节点在子节点之前”各有不同的合规结果。用错了场景，后面的判定也会错。
  \item 以合规初始结果和样例元素为起点，最多另用一个经证实不在样例中的同类型元素。其他字段固定，每份结果至多插入、删除、替换或对调两次。按改动次数、操作英文名字母序、位置和元素标识排序，含初始结果取前 64 份不同的原始返回值，保留重复项与顺序。判断前用案例和结果的内容指纹确定编号与展示次序。只有其他相关条款都满足，才把错误归给当前条款。把同一结果交给原测试的接收位置，照常解析并执行断言。违规结果通过算漏检，合规结果在指定断言失败算误拒。条款含糊、无法忠实替换返回值或其他位置先失败，将对应判断记为未知。另存依据、改动步骤和失败行。用 python-docx 1.2.0 读出的真实对象重放，原断言接受了 1 种正确顺序和全部 5 种错序。
\noindent\emph{A 是原断言的判断，L 是需求对同一结果的判断，1 表示接受。D 为 +1 表示放过违规，\(- 1\) 表示错拒合规，0 表示一致。有一项无法判断就记为未知。}
\par\noindent{\small\ttfamily [eq~1, raw LaTeX, not typeset] D=A-L,\textbackslash\{\}qquad A,L\textbackslash\{\}in\textbackslash\{\}\{0,1\textbackslash\{\}\}}
\par\emph{为什么：} 一份具体结果就能让别人复核原断言在哪里放过了违规，或错拒了合规。产品在相同输入下会不会产生这个结果，还要另找证据。
\end{enumerate}
\subsection*{选择一处断言修改}
\emph{检验指出具体哪一条判错，是否有助于选对改法。}
\begin{enumerate}[leftmargin=*,start=4]
  \item 两组都看到需求、样例、原测试、全部结果和原判定。按原判定、返回类型、字段、长度及改动操作序列分组。组内将判断一致和不一致的记录各按完整结果的标识排序，首行配首行，第二行配第二行。第一组保留原差值，第二组交换每对的两个差值。原判定减去分给该行的差值，就是新断言应给的答案。编辑器和筛选程序只用本组分配结果。隐藏原标签、理由、初始样例标记、配对和按真实错判挑出的例子。未知或未配对的行照常展示，不计分。一次最多提出 4 种只改一个条件的写法。合格改法须到达指定断言，并给对全部指定答案。按语法树最少增删改几个节点挑选，同分依次按源码位置和代码文本排序。无指定答案或无合格改法就保留原测试。锁定修改后再公布每种实际错判的首例、对应初始结果和配对数量。
\noindent\emph{A 是原断言的判断，D* 是分给这条结果的差值，两者算出修改后的断言应给出的判断 Y。只在 A 相同的结果之间对调差值，Y 仍为 0 或 1。}
\par\noindent{\small\ttfamily [eq~2, raw LaTeX, not typeset] Y=A-D\textasciicircum{}\{\textbackslash\{\}ast\}}
\par\emph{为什么：} 两边都看见相同的例子，也收到相同的一组差值。只对调差值归属，就能检验认准具体判错的那一条有没有帮助。
\end{enumerate}
\subsection*{核对诊断，检验真实缺陷}
\emph{先确认资料足以支持哪些判断，再把选定的改法放到独立实现上检验。}
\begin{enumerate}[leftmargin=*,start=5]
  \item 先保存仅凭公开资料得到的诊断。在另一份副本上，用相同模型、时间和样本上限，加入开发阶段可用的实现运行记录再做一次。两位读者在不知道组别和原断言结论的情况下，各自核对哪些结果符合需求。另行重放原断言。分别以全部筛过的案例和符合条件的案例为分母，列出确认、有争议、未知和缺资料的数量。每份开发运行记录都注明版本、输入样例和原始返回值。【作者需决定：比较两组之前，先定哪些开发版本及哪些运行可提供记录。】两人各自引用原文，给出允许、禁止或未知。两人都支持所记标签，且独立重放得到同样的原判定，才确认该错判。一个案例无论有几处错判都只计一次，另列仅加入运行记录后能确认的案例。没发现错判的案例仍在相应分母内。
\par\emph{为什么：} 这能量出仅靠公开资料可以确认多少诊断，加上运行记录又能补出多少。
  \item 比较原测试、盲改、直接把需求写成断言、无引导的同类型结果改动，以及两种差值反馈。直接翻译需求的一方也能看到完整需求、样例、动作、读取结果的代码、原断言和同一批无标签结果。它能在相同的调用和重放次数内运行自己写的判定，最后选一个条件修改。各组记录调用次数、重放次数和人工核对时间，放弃修改的案例以原测试参加评估。盲改组用公开材料复查原断言。直接翻译组把需求写成判断，例如逐个核对两个预期文件名各出现几次。无引导改动只按类型选择元素和位置。保存各组实收的输入和最终改法，固定样例、动作和读取结果的代码。
\par\emph{为什么：} 直接翻译已经用到了需求本身。如果它在同等预算下得出相同的检错与误拒结果，现有证据就不足以说明重放带来了额外收益。
  \item 评分前，请未看过反例和改法的评估者登记历史缺陷、对应修复版、需求原句、样例和运行命令。对合规变体，记下哪个独立实现在相同输入下确实产生了这种允许的变化。【作者需决定：先选定独立档案来源，并约定哪些发布版本或独立实现可算作不同的合规变体。】整条测试在错误版的指定断言失败、在它自己的修复版通过，才算检出。同一缺陷每个项目只计一次。每个项目用新增检出数减去原能检出而现在漏掉的数量，除以全部可评估缺陷数，再对项目求平均。合规变体另算，以全部可评估变体为分母，数出原先通过、如今在该断言失败的数量。单列无关失败、无法运行及已有条款依据的基准变异。哪一项缺样本，就标明该项未验证。现有 833 对测试与未检出变异记录中，只有 3 对明确关联到了条款。
\noindent\emph{项目 p 有 \(B_{p}\) 个可评估的缺陷，\(G_{p}\) 个新增检出，\(M_{p}\) 个原能检出、修改后漏掉。\(V_{p}\) 个可评估的合规变体中，\(J_{p}\) 个原先通过、修改后被拒绝。哪一项分母为零，哪一项就暂时无法评估。}
\par\noindent{\small\ttfamily [eq~3, raw LaTeX, not typeset] g\_p=\textbackslash\{\}frac\{G\_p-M\_p\}\{B\_p\},\textbackslash\{\}qquad r\_p=\textbackslash\{\}frac\{J\_p\}\{V\_p\}}
\par\emph{为什么：} 缺陷版和修复版都失败的测试得 0 次检出。检查实际出现过的合法变体，才能知道这次修改会不会把允许的行为也挡住。
\end{enumerate}
\end{document}
